\scanpage{11}\mspage{5}
\[
\left\{\begin{aligned}
&\text{l'image inverse de } \mathbb{F}_2\times\{0\} &&\text{est }\deleted{\unsure{\mathfrak{D}|\mathscr{D}}}V_S\\
&\rule{5em}{0.3pt}\quad \{0\}\times\mathbb{F}_2 &&\text{est }\deleted{\mathfrak{D}}V_A\\
&\text{donc}\quad\rule{5em}{0.3pt}\quad\substack{\text{diagonale}\\\text{de }\mathbb{F}_2\times\mathbb{F}_2} &&\text{est }\mathfrak{D}^+
\end{aligned}\right.
\]
ou encore \deleted{\unsure{l'i}}

\[
\left\{\begin{aligned}
&\text{l'image inverse de }(1,0) &&\text{est }C\subset\mathfrak{D}^{\unsure{4|*|\bullet}}\\
&\rule{5em}{0.3pt}\quad(0,1) &&\text{est }D\subset\mathfrak{D}^{\unsure{4|*}}\\
&\rule{5em}{0.3pt}\quad(1,1) &&\text{est }\underline\omega_Q\subset\mathfrak{D}^{\unsure{4|*}}
\end{aligned}\right.
\]
Les sous-groupes d'\deleted{ordre deux} de $\mathfrak{D}$ correspondent aux éléments d'ordre $2$, il y a comme tels éléments

a) l'antipodisme $\underline a$, engendrant le centre $\{1,\underline a\}$

b) Les deux éléments de $C=V_S\setminus\{1,a\}$, engendrant resp. les deux droites-coordonnées de $V_S\simeq\mathbb{F}_2^C$

c) Les deux éléments de $D=V_A\setminus\{1,a\}$, engendrant resp les deux droites coordonnées de $V_A\simeq\mathbb{F}_2^D$.
% slanted
\margin{NB. ces quatre éléments n'étant pas centraux, les sous-groupes d'ordre $2$ correspondants ne sont pas distingués ; leurs normalisateurs sont resp. $V_S$ et $V_A$. Donc parmi les huit ss-groupes de $\mathfrak{D}$ $\ne\{1\},\mathfrak{D}$, il y en a quatre qui sont distingués ($V_S,V_A,\mathfrak{D}^+,\{1,\underline a\}$) et quatre qui \unsure{le|ne le} sont pas.}

(NB les \deleted{\unsure{autres}} $3$ éléments de $\mathfrak{D}^{\unsure{4|*}}$ sont d'ordre $1$, ou d'ordre $4$ :
\[
\underbrace{\mathfrak{D}\simeq\{1\}\amalg\{a\}\amalg C\amalg D\amalg\underline\omega_Q}_{\text{déc. de }\mathfrak{D}\text{ en cinq orbites sous }\mathrm{int}(\mathfrak D)}\qquad)
\]
On a donc fait la liste complète des sous-groupes de $\mathfrak{D}$ différents de $\{1\}$ et de $\mathfrak{D}$ lui-$=$, il y en a trois d'ordre $4$ i.e. d'indice $2$, savoir $V_S,V_A$ et $\mathfrak{D}^+$, et cinq d'ordre $2$ savoir le centre $\{1,\underline a\}$, et les $4$ \inserted{$2$ paquets de $2$} sous-groupes d'ordre $2$ distincts de ce dernier, contenus respectivement dans $V_S$ et dans $V_A$.

On a associé à un carré $Q=(S,A,R)$ trois ensembles remarquables de cardinal $2$, savoir
\[
\unsure{D|B},S,\underline\omega_Q
\]
\margin{\unsure{4|$\{$}}
% archive sheet number: 10
